\section{Finite examples}
\label{sec:examples}

For a fixed target exponent, chain schemes much smaller than those of
Theorem~\ref{thm:main} suffice. This section lists certified examples
and compares them with the examples of~\citep[Section~10]{FinCon26}. No
other result depends on this section.

\subsection{Certified witnesses}
\label{subsec:examples-witnesses}

A \emph{witness} is given by a target \(\Phi_0\), a terminal pair
\((1,t)\), arities \(d_1,\ldots,d_J\ge2\) and integers \(N,k\ge1\). Its
chain is built backwards from \((1,t)\) by smallest parents
(Lemma~\ref{lem:smallest-parent}), which fixes the source pair and makes
every stage legal. Stage \(i\) uses a type of width \(N\) with the
largest sector count \(M(N,d_i)\) (Section~\ref{subsec:chains-potential}).
Proposition~\ref{prop:chain-counts} gives a restriction
\(r_*\MM_A\restr m_*\MM_B\) over \(\Q\) with the
counts~\eqref{eq:stage-chain-counts}, and we certify \(\Phi<\Phi_0\) for its
fixed-batch bound \(\Phi\)~\eqref{eq:chains-phi}, so that also
\(p<\Phi_0\). For example, the first row of Table~\ref{tab:witnesses}
is the single stage \((128,66)\to(1,128)\) of arity \(2\), legal since
\(4\le2\cdot66-128+1\).

\begin{proposition}[Witnesses]\label{prop:witnesses}
For each row of Table~\ref{tab:witnesses}, the restriction defined above
is a relative scheme over \(\Q\), its fixed-batch bound \(\Phi\) is less
than the entry in the column \(\Phi<\), hence less than \(\Phi_0\), and
\(\Phi-p\) and the numbers of decimal digits of \(A\) and \(B\) are as
stated.
\end{proposition}

\begin{proof}
The proof is computer-assisted, by the supplementary script
\path|shifted_chain_certificates.py|, which reproduces its report
\path|shifted-chain-certificates.json| byte for byte. The script checks
every stage against~\eqref{eq:stage-conditions}, and the maximality of each
type \(\tau=(\tau_1,\tau_2,\tau_3)=(\nK,\nO,\nmid)\) by the
inequalities \(w_i\tau_j\le w_j(\tau_i+1)\) for \(i\ne j\) with
\(\tau_j>0\), where \(w=(d-1,2,1)\), which say that moving one letter
from class \(j\) to class \(i\) does not increase the sector count;
these characterize the maximizers,
because the logarithm of~\eqref{eq:stage-sector-count} is \(\log N!\)
plus a sum of concave functions
\(\tau_i\mapsto\tau_i\log w_i-\log\tau_i!\). It forms the integers
\(d_i\), \(M_i\), \(2M_i-1\), \(E_i\), \(s\) and \(t\) exactly, encloses
their logarithms in intervals with dyadic endpoints (an
\(\operatorname{artanh}\) series with directed rounding and a tail
bound), and evaluates \(\Phi\), \(p\), \(\Phi-p\), \(\log_{10}A\) and
\(\log_{10}B\) in exact interval arithmetic. If the enclosure of
\(\log_{10}A\) or \(\log_{10}B\) contains an integer, the digit count is
decided exactly. Printed bounds are rounded outwards.
\end{proof}

\begin{table}[t]
\centering
\small
\setlength{\tabcolsep}{3pt}
\begin{tabular}{lrlrrrlrrr}
\toprule
\(\Phi_0\)&\(J\)&arities&\(t\)&\(N\)&\(k\)&\(\Phi<\)&\(\Phi-p\)
 &digits of \(A\)&digits of \(B\)\\
\midrule
\(3\)&\(1\)&\(2\)&\(128\)&\(24\)&\(2\)&\(2.999745128438\)
 &\(1.34\cdot10^{-3}\)&\(29\)&\(255\)\\
\(2.77429\)&\(1\)&\(2\)&\(5\)&\(215\)&\(10\)&\(2.774233389609\)
 &\(7.08\cdot10^{-5}\)&\(1\,295\)&\(5\,547\)\\
\(2.7\)&\(5\)&\(2\)&\(3\)&\(147\)&\(8\)&\(2.699881286935\)
 &\(6.71\cdot10^{-5}\)&\(3\,541\)&\(8\,025\)\\
\(2.5\)&\(15\)&\(2\)&\(3\)&\(264\)&\(14\)&\(2.499903515815\)
 &\(8.39\cdot10^{-6}\)&\(33\,379\)&\(69\,247\)\\
\(2.3\)&\(125\)&\(2\), ramp \(3\)&\(1\)&\(1\,999\)&\(80\)
 &\(2.299993369880\)&\(2.54\cdot10^{-8}\)&\(12\,371\,380\)
 &\(24\,208\,459\)\\
\(2.26\)&\(1\,264\)&\(4\), ramp \(11\)&\(1\)&\(28\,336\)&\(1\,118\)
 &\(2.259999981185\)&\(2.16\cdot10^{-11}\)&\(48\,859\,824\,830\)
 &\(62\,827\,651\,353\)\\
\midrule
\(2.84063\)&\(1\)&\(2\)&\(5\)&\(93\)&\(5\)&\(2.840553213477\)
 &\(3.31\cdot10^{-4}\)&\(280\)&\(1\,190\)\\
\(2.67283\)&\(7\)&\(2\)&\(3\)&\(138\)&\(8\)&\(2.672186713007\)
 &\(5.51\cdot10^{-5}\)&\(4\,653\)&\(10\,114\)\\
\bottomrule
\end{tabular}
\caption{Certified chain schemes (Proposition~\ref{prop:witnesses}),
built backwards from \((1,t)\) by smallest parents. Column \(J\): the
number of stages. Column ``arities'': the arity \(d\) of every stage,
or ``\(d\), ramp~\(R\)'' for \(J-2(R-1)\) stages of arity \(d\) followed
by the \(2(R-1)\) stages of Lemma~\ref{lem:ramp} from \((R,R)\) to
\((1,1)\), of arities \(3j\) and \(3j-1\) for \(j=R-1,\ldots,1\), in
this order from the source. Column \(\Phi<\): certified upper bound for
\(\Phi\). Column \(\Phi-p\): the value of \(\log2/\log(B/A)\), rounded to
three significant digits. The last two rows are the comparison points of
Section~\ref{subsec:examples-comparison}.}
\label{tab:witnesses}
\end{table}

The source pairs are \((128,66)\), \((5,4)\), \((18,27)\),
\((3098,5205)\), a pair of \(28\)-digit integers, a pair of
\(166\)-digit integers, \((5,4)\) and \((49,78)\); the report lists them
with the types and intervals for \(p\). The rows are selected, not
optimized; we do not claim that any of them is the smallest chain scheme
below its target. The \(2.3\) and \(2.26\) rows have the ramp-and-bulk
shape of Theorem~\ref{thm:main}, but much smaller parameters than that
theorem uses. The column \(\Phi-p\) is below \(1.4\cdot10^{-3}\) in
every row and below \(10^{-5}\) for the targets \(2.5\), \(2.3\)
and~\(2.26\).

\subsection{Comparison with the preceding paper}
\label{subsec:examples-comparison}

Table~2 of~\citep{FinCon26} bounds the same quantity \(\Phi\)
\citep[(10.1)]{FinCon26} for three schemes of the construction there,
with target indices \(a=4\), \(8\) and \(12\) in the notation
of~\citep{FinCon26}. For the last two it gives the certified upper
bounds \(2.848599165\) and \(2.675820359\), the floating-point values
\(2.8406393\) and \(2.6728401\) of \(\Phi\), and source sides with \(11\,877\,489\,885\) and \(234\,905\,651\,862\) digits.
The last two rows of Table~\ref{tab:witnesses} have targets about
\(10^{-5}\) below these floating-point values. They are the one-stage
chain \((5,4)\to(1,5)\) of arity \(2\), with \(A=2^{930}\), and a chain
of seven stages of arity \(2\) from \((49,78)\) to \((1,3)\), with
\(A=2^{15456}\). Their source sides have \(280\) and \(4\,653\) digits,
fewer by factors of about \(4.2\cdot10^7\) and \(5.0\cdot10^7\), and
even their target sides, with \(1\,190\) and \(10\,114\) digits, have
fewer digits than the source sides there. Pooling lowers the \(\Phi\) of
those schemes by less than \(10^{-7}\)~\citep[Section~10.2]{FinCon26},
so the targets are also below their pooled exponents, provided the
floating-point values are accurate to \(10^{-6}\).

Near \(9/4\) the gap is much larger. By
\citep[Corollary~11.2(a)]{FinCon26}, every scheme of the construction
there with \(p<3\) has \(p>\frac94+9\log(6/5)/(2\log(2\log A))\). For
\(p<2.26\) this forces \(2\log A>(6/5)^{450}\), so \(A\) has more than
\((6/5)^{450}/(2\log10)>9.29\cdot10^{34}\) digits, against
\(48\,859\,824\,830\) for the \(2.26\) row. For \(p<2.3\) it forces only
\(2\log A>(6/5)^{90}\), that is, more than \(2.9\cdot10^6\) digits, so it
does not separate the two constructions there.

These comparisons concern the schemes as constructed, with the source
supplying every auxiliary. Forward recycling and composition with other
relative schemes, for example with the classical algorithm
\(n^3\MM_1\restr\MM_n\) of source side \(1\), change \(A\) and \(B\)
(Section~\ref{subsec:limits-class}).

Far from \(9/4\), chain schemes are not the smallest schemes known.
Strassen's algorithm~\citep{Str69} is the relative scheme
\(7\MM_1\restr\MM_2\), with \(A=1\), \(B=2\) and \(p=\log_27<2.808\),
below the targets \(3\) and \(2.84063\). Schwartz and Zwecher~\citep{SZ26} report an algorithm
with exponent \(2.773203\) whose base case has side below \(1000\),
below the target \(2.77429\). The advantage of chain schemes lies near
\(9/4\), where, to our knowledge, the only other explicit sizes are
those of~\citep{FinCon26}, which grow doubly exponentially in
\(1/\eps\)~\citep[Corollary~11.2(b)]{FinCon26}.
